\documentclass[12pt]{article}
\usepackage{amsmath}
\usepackage{geometry}
\geometry{a4paper, margin=1in}

\title{Complete Lagrangian Derivation for the Double Inverted Pendulum on a Cart (DIPC)}
\author{Phillipp Gery}
\date{\today}

\begin{document}
	
	\maketitle
	
	\section{Derivation of the Lagrangian}
	The Lagrangian is defined as $L = T - V$, where $T$ is the total kinetic energy and $V$ is the total potential energy of the system.
	
	\subsection{Kinetic Energy (T)}
	The total kinetic energy is the sum of the kinetic energies of the cart ($T_M$), the first mass ($T_{m1}$), and the second mass ($T_{m2}$).
	
	\subsubsection{Cart Kinetic Energy ($T_M$)}
	The cart moves only in the x-direction.
	\begin{align*}
		\text{Position:} \quad & (x, 0) \\
		\text{Velocity:} \quad & (\dot{x}, 0) \\
		T_M &= \frac{1}{2} M \dot{x}^2
	\end{align*}
	
	\subsubsection{Mass 1 Kinetic Energy ($T_{m1}$)}
	First, we define the absolute position of $m_1$ and differentiate it to find its velocity.
	\begin{align*}
		x_1 &= x + l_1 \sin(\theta_1) \\
		y_1 &= l_1 \cos(\theta_1)
	\end{align*}
	The velocity components are:
	\begin{align*}
		\dot{x}_1 &= \dot{x} + l_1 \dot{\theta}_1 \cos(\theta_1) \\
		\dot{y}_1 &= -l_1 \dot{\theta}_1 \sin(\theta_1)
	\end{align*}
	The kinetic energy is $T_{m1} = \frac{1}{2} m_1 (\dot{x}_1^2 + \dot{y}_1^2)$:
	\begin{align*}
		T_{m1} &= \frac{1}{2} m_1 \left[ (\dot{x} + l_1 \dot{\theta}_1 \cos(\theta_1))^2 + (-l_1 \dot{\theta}_1 \sin(\theta_1))^2 \right] \\
		&= \frac{1}{2} m_1 \left[ \dot{x}^2 + 2l_1\dot{x}\dot{\theta}_1\cos(\theta_1) + l_1^2\dot{\theta}_1^2\cos^2(\theta_1) + l_1^2\dot{\theta}_1^2\sin^2(\theta_1) \right] \\
		&= \frac{1}{2} m_1 \left[ \dot{x}^2 + 2l_1\dot{x}\dot{\theta}_1\cos(\theta_1) + l_1^2\dot{\theta}_1^2(\cos^2(\theta_1) + \sin^2(\theta_1)) \right] \\
		T_{m1} &= \frac{1}{2} m_1 (\dot{x}^2 + l_1^2 \dot{\theta}_1^2 + 2 l_1 \dot{x} \dot{\theta}_1 \cos(\theta_1))
	\end{align*}
	
	\subsubsection{Mass 2 Kinetic Energy ($T_{m2}$)}
	The position of $m_2$ depends on the positions of the cart and $m_1$.
	\begin{align*}
		x_2 &= x + l_1 \sin(\theta_1) + l_2 \sin(\theta_2) \\
		y_2 &= l_1 \cos(\theta_1) + l_2 \cos(\theta_2)
	\end{align*}
	The velocity components are:
	\begin{align*}
		\dot{x}_2 &= \dot{x} + l_1 \dot{\theta}_1 \cos(\theta_1) + l_2 \dot{\theta}_2 \cos(\theta_2) \\
		\dot{y}_2 &= -l_1 \dot{\theta}_1 \sin(\theta_1) - l_2 \dot{\theta}_2 \sin(\theta_2)
	\end{align*}
	The kinetic energy $T_{m2} = \frac{1}{2} m_2 (\dot{x}_2^2 + \dot{y}_2^2)$ is found by summing the squares of the velocity components. This yields:
	\begin{align*}
		T_{m2} = \frac{1}{2} m_2 \Big[ & \dot{x}^2 + l_1^2 \dot{\theta}_1^2 + l_2^2 \dot{\theta}_2^2 + 2\dot{x}l_1\dot{\theta}_1\cos(\theta_1) + 2\dot{x}l_2\dot{\theta}_2\cos(\theta_2) \\
		&+ 2l_1l_2\dot{\theta}_1\dot{\theta}_2\cos(\theta_1-\theta_2) \Big]
	\end{align*}
	
	\subsubsection{Total Kinetic Energy (T)}
	Summing the components $T = T_M + T_{m1} + T_{m2}$ and grouping terms gives the total kinetic energy:
	\begin{align*}
		T = & \frac{1}{2} (M + m_1 + m_2) \dot{x}^2 + \frac{1}{2} (m_1 + m_2) l_1^2 \dot{\theta}_1^2 + \frac{1}{2} m_2 l_2^2 \dot{\theta}_2^2 \\
		&+ (m_1 + m_2) l_1 \dot{x} \dot{\theta}_1 \cos(\theta_1) + m_2 l_2 \dot{x} \dot{\theta}_2 \cos(\theta_2) \\
		&+ m_2 l_1 l_2 \dot{\theta}_1 \dot{\theta}_2 \cos(\theta_1 - \theta_2)
	\end{align*}
	
	\subsection{Potential Energy (V)}
	The potential energy is determined by the height of each mass relative to the cart ($y=0$).
	\begin{align*}
		V_M &= 0 \\
		V_{m1} &= m_1 g y_1 = m_1 g l_1 \cos(\theta_1) \\
		V_{m2} &= m_2 g y_2 = m_2 g (l_1 \cos(\theta_1) + l_2 \cos(\theta_2))
	\end{align*}
	The total potential energy $V = V_M + V_{m1} + V_{m2}$ is:
	\begin{equation*}
		V = (m_1 + m_2) g l_1 \cos(\theta_1) + m_2 g l_2 \cos(\theta_2)
	\end{equation*}
	
	\subsection{Assembling the Lagrangian}
	Substituting the final expressions for $T$ and $V$ into $L=T-V$ gives the full Lagrangian for the system:
	\begin{align*}
		L &= \frac{1}{2} (M + m_1 + m_2) \dot{x}^2 + \frac{1}{2} (m_1 + m_2) l_1^2 \dot{\theta}_1^2 + \frac{1}{2} m_2 l_2^2 \dot{\theta}_2^2 \\
		&\quad + (m_1 + m_2) l_1 \dot{x} \dot{\theta}_1 \cos(\theta_1) + m_2 l_2 \dot{x} \dot{\theta}_2 \cos(\theta_2) \\
		&\quad + m_2 l_1 l_2 \dot{\theta}_1 \dot{\theta}_2 \cos(\theta_1 - \theta_2) - (m_1 + m_2) g l_1 \cos(\theta_1) - m_2 g l_2 \cos(\theta_2)
	\end{align*}
	
	\section{Derivation of the Equations of Motion}
	The Lagrange equations of motion are derived from the formula:
	\begin{equation}
		\frac{d}{dt} \left( \frac{\partial L}{\partial \dot{q}_i} \right) - \frac{\partial L}{\partial q_i} = Q_i
	\end{equation}
	where $L$ is the Lagrangian, $q_i$ are the generalized coordinates, and $Q_i$ are the generalized non-conservative forces. For this system, the coordinates are $q = [x, \theta_1, \theta_2]^T$ and the forces are $Q_x = u$, $Q_{\theta_1} = 0$, and $Q_{\theta_2} = 0$.
	
	\subsection{Equation for Coordinate $x$}
	
	\subsubsection*{Step 1: Calculate $\frac{\partial L}{\partial \dot{x}}$}
	\[
	\frac{\partial L}{\partial \dot{x}} = (M + m_1 + m_2) \dot{x} + (m_1 + m_2) l_1 \dot{\theta}_1 \cos(\theta_1) + m_2 l_2 \dot{\theta}_2 \cos(\theta_2)
	\]
	
	\subsubsection*{Step 2: Calculate $\frac{d}{dt} \left( \frac{\partial L}{\partial \dot{x}} \right)$}
	We apply the total time derivative, using the product and chain rules.
	\begin{align*}
		\frac{d}{dt} \left( \frac{\partial L}{\partial \dot{x}} \right) &= (M + m_1 + m_2) \ddot{x} + (m_1 + m_2) l_1 \left( \ddot{\theta}_1 \cos(\theta_1) - \dot{\theta}_1^2 \sin(\theta_1) \right) \\
		&\quad + m_2 l_2 \left( \ddot{\theta}_2 \cos(\theta_2) - \dot{\theta}_2^2 \sin(\theta_2) \right)
	\end{align*}
	
	\subsubsection*{Step 3: Calculate $\frac{\partial L}{\partial x}$}
	The Lagrangian $L$ does not explicitly depend on $x$.
	\[
	\frac{\partial L}{\partial x} = 0
	\]
	
	\subsubsection*{Step 4: Assemble the Equation}
	Using $\frac{d}{dt} \left( \frac{\partial L}{\partial \dot{x}} \right) - \frac{\partial L}{\partial x} = u$, the final equation of motion for $x$ is:
	\begin{equation}
		(M \!+\! m_1 \!+\! m_2) \ddot{x} + (m_1 \!+\! m_2) l_1 \ddot{\theta}_1 \cos(\theta_1) + m_2 l_2 \ddot{\theta}_2 \cos(\theta_2) - (m_1 \!+\! m_2) l_1 \dot{\theta}_1^2 \sin(\theta_1) - m_2 l_2 \dot{\theta}_2^2 \sin(\theta_2) = u
	\end{equation}
	
	\subsection{Equation for Coordinate $\theta_1$}
	
	\subsubsection*{Step 1: Calculate $\frac{\partial L}{\partial \dot{\theta}_1}$}
	\[
	\frac{\partial L}{\partial \dot{\theta}_1} = (m_1 + m_2) l_1^2 \dot{\theta}_1 + (m_1 + m_2) l_1 \dot{x} \cos(\theta_1) + m_2 l_1 l_2 \dot{\theta}_2 \cos(\theta_1 - \theta_2)
	\]
	
	\subsubsection*{Step 2: Calculate $\frac{d}{dt} \left( \frac{\partial L}{\partial \dot{\theta}_1} \right)$}
	\begin{align*}
		\frac{d}{dt} \left( \frac{\partial L}{\partial \dot{\theta}_1} \right) &= (m_1 + m_2) l_1^2 \ddot{\theta}_1 + (m_1 + m_2) l_1 \left( \ddot{x}\cos(\theta_1) - \dot{x}\dot{\theta}_1\sin(\theta_1) \right) \\
		&\quad + m_2 l_1 l_2 \left( \ddot{\theta}_2\cos(\theta_1 - \theta_2) - \dot{\theta}_2(\dot{\theta}_1 - \dot{\theta}_2)\sin(\theta_1 - \theta_2) \right)
	\end{align*}
	
	\subsubsection*{Step 3: Calculate $\frac{\partial L}{\partial \theta_1}$}
	\[
	\frac{\partial L}{\partial \theta_1} = -(m_1 + m_2) l_1 \dot{x} \dot{\theta}_1 \sin(\theta_1) - m_2 l_1 l_2 \dot{\theta}_1 \dot{\theta}_2 \sin(\theta_1 - \theta_2) + (m_1 + m_2) g l_1 \sin(\theta_1)
	\]
	
	\subsubsection*{Step 4: Assemble the Equation}
	Using $\frac{d}{dt} \left( \frac{\partial L}{\partial \dot{\theta}_1} \right) - \frac{\partial L}{\partial \theta_1} = 0$, several terms cancel, leaving the final equation of motion for $\theta_1$:
	\begin{equation}
		(m_1 \!+\! m_2) l_1 \ddot{x} \cos(\theta_1) + (m_1 \!+\! m_2) l_1^2 \ddot{\theta}_1 + m_2 l_1 l_2 \ddot{\theta}_2 \cos(\theta_1 \!-\! \theta_2) + m_2 l_1 l_2 \dot{\theta}_2^2 \sin(\theta_1 \!-\! \theta_2) - (m_1 \!+\! m_2) g l_1 \sin(\theta_1) = 0
	\end{equation}
	
	\subsection{Equation for Coordinate $\theta_2$}
	
	\subsubsection*{Step 1: Calculate $\frac{\partial L}{\partial \dot{\theta}_2}$}
	\[
	\frac{\partial L}{\partial \dot{\theta}_2} = m_2 l_2^2 \dot{\theta}_2 + m_2 l_2 \dot{x} \cos(\theta_2) + m_2 l_1 l_2 \dot{\theta}_1 \cos(\theta_1 - \theta_2)
	\]
	
	\subsubsection*{Step 2: Calculate $\frac{d}{dt} \left( \frac{\partial L}{\partial \dot{\theta}_2} \right)$}
	\begin{align*}
		\frac{d}{dt} \left( \frac{\partial L}{\partial \dot{\theta}_2} \right) &= m_2 l_2^2 \ddot{\theta}_2 + m_2 l_2 \left( \ddot{x}\cos(\theta_2) - \dot{x}\dot{\theta}_2\sin(\theta_2) \right) \\
		&\quad + m_2 l_1 l_2 \left( \ddot{\theta}_1\cos(\theta_1 - \theta_2) - \dot{\theta}_1(\dot{\theta}_1 - \dot{\theta}_2)\sin(\theta_1 - \theta_2) \right)
	\end{align*}
	
	\subsubsection*{Step 3: Calculate $\frac{\partial L}{\partial \theta_2}$}
	\[
	\frac{\partial L}{\partial \theta_2} = -m_2 l_2 \dot{x} \dot{\theta}_2 \sin(\theta_2) + m_2 l_1 l_2 \dot{\theta}_1 \dot{\theta}_2 \sin(\theta_1 - \theta_2) + m_2 g l_2 \sin(\theta_2)
	\]
	
	\subsubsection*{Step 4: Assemble the Equation}
	Using $\frac{d}{dt} \left( \frac{\partial L}{\partial \dot{\theta}_2} \right) - \frac{\partial L}{\partial \theta_2} = 0$, terms again cancel, yielding the final equation of motion for $\theta_2$:
	\begin{equation}
		m_2 l_2 \ddot{x} \cos(\theta_2) + m_2 l_1 l_2 \ddot{\theta}_1 \cos(\theta_1 \!-\! \theta_2) + m_2 l_2^2 \ddot{\theta}_2 - m_2 l_1 l_2 \dot{\theta}_1^2 \sin(\theta_1 \!-\! \theta_2) - m_2 g l_2 \sin(\theta_2) = 0
	\end{equation}
	
	\newpage
	\section{Final Overview and Matrix Form}
	
	This section summarizes the three complete, non-linear equations of motion and presents them in the compact matrix form common in robotics and control theory.
	
	\subsection{Final Equations of Motion}
	The three derived equations are:
	
	\begin{enumerate}
		\item \textbf{For coordinate $x$:}
		\begin{align*}
			(M \!+\! m_1 \!+\! m_2) \ddot{x} + (m_1 \!+\! m_2) l_1 \ddot{\theta}_1 \cos(\theta_1) + m_2 l_2 \ddot{\theta}_2 \cos(\theta_2) \\ 
			- (m_1 \!+\! m_2) l_1 \dot{\theta}_1^2 \sin(\theta_1) - m_2 l_2 \dot{\theta}_2^2 \sin(\theta_2) = u
		\end{align*}
		
		\item \textbf{For coordinate $\theta_1$:}
		\begin{align*}
			(m_1 \!+\! m_2) l_1 \ddot{x} \cos(\theta_1) + (m_1 \!+\! m_2) l_1^2 \ddot{\theta}_1 + m_2 l_1 l_2 \ddot{\theta}_2 \cos(\theta_1 \!-\! \theta_2) \\
			+ m_2 l_1 l_2 \dot{\theta}_2^2 \sin(\theta_1 \!-\! \theta_2) - (m_1 \!+\! m_2) g l_1 \sin(\theta_1) = 0
		\end{align*}
		
		\item \textbf{For coordinate $\theta_2$:}
		\begin{align*}
			m_2 l_2 \ddot{x} \cos(\theta_2) + m_2 l_1 l_2 \ddot{\theta}_1 \cos(\theta_1 \!-\! \theta_2) + m_2 l_2^2 \ddot{\theta}_2 \\
			- m_2 l_1 l_2 \dot{\theta}_1^2 \sin(\theta_1 \!-\! \theta_2) - m_2 g l_2 \sin(\theta_2) = 0
		\end{align*}
	\end{enumerate}
	
	\subsection{State-Space Matrix Form}
	
	The state vector for the system is defined as:
	\begin{equation}
		\mathbf{x} = \begin{pmatrix} x_1 \\ x_2 \\ x_3 \\ x_4 \\ x_5 \\ x_6 \end{pmatrix} = \begin{pmatrix} x \\ \theta_1 \\ \theta_2 \\ \dot{x} \\ \dot{\theta}_1 \\ \dot{\theta}_2 \end{pmatrix}
	\end{equation}
	
	Non-Linear State Equation:\\
	\begin{equation}
		\dot{\mathbf{x}} = f(\mathbf{x}, u)
	\end{equation}\\


	The first three components of the state-derivative vector are kinematic definitions derived directly from the state vector itself:
		\begin{align*}
			\dot{x}_1 &= \dot{x} = x_4 \\
			\dot{x}_2 &= \dot{\theta}_1 = x_5 \\
			\dot{x}_3 &= \dot{\theta}_2 = x_6
		\end{align*}




	The last three components (the accelerations) are found by solving the system's dynamic EoM from Lagrangian for each elements.\\
	This part will be continued in MATLAB.
	
	
%	 These equations can be written in the more compact and standard state-space matrix form for robotic manipulators as done from Alexander Bogdanov in "Optimal Control of a Double Inverted Pendulum on a Cart":
%	\begin{equation}
%		\mathbf{D}(q)\ddot{q} + \mathbf{C}(q, \dot{q})\dot{q} + \mathbf{G}(q) = \mathbf{H}u
%	\end{equation}
%	where the generalized coordinate vector is $q = [x, \theta_1, \theta_2]^T$. This results in the following for our case:
%	\begin{equation}
%		\mathbf{D}(\mathbf{x}) \begin{pmatrix} \ddot{x} \\ \ddot{\theta}_1 \\ \ddot{\theta}_2 \end{pmatrix} + \mathbf{C}(\mathbf{x}) \begin{pmatrix} \dot{x} \\ \dot{\theta}_1 \\ \dot{\theta}_2 \end{pmatrix} + \mathbf{G}(\mathbf{x}) = \mathbf{H}u
%	\end{equation}
%	Solving for the acceleration vector gives:
%	\begin{equation}
%		\begin{pmatrix} \dot{x}_4 \\ \dot{x}_5 \\ \dot{x}_6 \end{pmatrix} = \begin{pmatrix} \ddot{x} \\ \ddot{\theta}_1 \\ \ddot{\theta}_2 \end{pmatrix} = \mathbf{D}(\mathbf{x})^{-1} \left( \mathbf{H}u - \mathbf{C}(\mathbf{x})\begin{pmatrix} x_4 \\ x_5 \\ x_6 \end{pmatrix} - \mathbf{G}(\mathbf{x}) \right)
%	\end{equation}
%	The matrices, expressed in terms of the state variables, are defined below.
%	
%
%	\[
%	\mathbf{D}(\mathbf{x}) = 
%	\begin{pmatrix}
%		M+m_1+m_2 & (m_1+m_2)l_1\cos(x_2) & m_2l_2\cos(x_3) \\
%		(m_1+m_2)l_1\cos(x_2) & (m_1+m_2)l_1^2 & m_2l_1l_2\cos(x_2-x_3) \\
%		m_2l_2\cos(x_3) & m_2l_1l_2\cos(x_2-x_3) & m_2l_2^2
%	\end{pmatrix}
%	\]
%	
%	\[
%	\mathbf{C}(\mathbf{x})\dot{\mathbf{x}} = 
%	\begin{pmatrix}
%		0&-(m_1+m_2)l_1x_5\sin(x_2) &- m_2l_2x_6\sin(x_3) \\
%		0&0&m_2l_1l_2x_6\sin(x_2-x_3) \\
%		0&-m_2l_1l_2x_5\sin(x_2-x_3)&0
%	\end{pmatrix}
%	\]
%	
%
%	\[
%	\mathbf{G}(\mathbf{x}) = 
%	\begin{pmatrix}
%		0 \\
%		(m_1+m_2)g l_1\sin(x_2) \\
%		m_2g l_2\sin(x_3)
%	\end{pmatrix}
%	\]
%	
%
%	\[
%	\mathbf{H} = 
%	\begin{pmatrix}
%		1 \\ 0 \\ 0
%	\end{pmatrix}
%	\]
	
	\subsubsection{Output Equation: $\mathbf{y} = \mathbf{C}\mathbf{x} + \mathbf{D}\mathbf{u}$}
	If we choose the outputs to be the positions, then the output vector is $\mathbf{y} = [x, \theta_1, \theta_2]^T$. In this  case, the relationship between the states and outputs is linear and can be written in the standard form.
	

	\[
	\mathbf{C} =
	\begin{pmatrix}
		1 & 0 & 0 & 0 & 0 & 0 \\
		0 & 1 & 0 & 0 & 0 & 0 \\
		0 & 0 & 1 & 0 & 0 & 0
	\end{pmatrix}
	\]
	

	Since the input force does not instantaneously affect the output positions, the feedthrough matrix is zero.
	\[
	\mathbf{D} =
	\begin{pmatrix}
		0 \\
		0 \\
		0
	\end{pmatrix}
	\]
	
\end{document}